
\documentclass[12pt]{article}
\usepackage{graphicx}
%%%%%%%%%%%%%%%%%%%%%%%%%%%%%%%%%%%%%%%%%%%%%%%%%%%%%%%%%%%%%%%%%%%%%%%%%%%%%%%%%%%%%%%%%%%%%%%%%%%%%%%%%%%%%%%%%%%%%%%%%%%%%%%%%%%%%%%%%%%%%%%%%%%%%%%%%%%%%%%%%%%%%%%%%%%%%%%%%%%%%%%%%%%%%%%%%%%%%%%%%%%%%%%%%%%%%%%%%%%%%%%%%%%%%
\usepackage{amsmath}
\usepackage[compat2]{geometry}

\setcounter{MaxMatrixCols}{10}
%TCIDATA{TCIstyle=LaTeX article (bright).cst}

%TCIDATA{OutputFilter=LATEX.DLL}
%TCIDATA{Version=5.50.0.2953}
%TCIDATA{<META NAME="SaveForMode" CONTENT="3">}
%TCIDATA{BibliographyScheme=Manual}
%TCIDATA{Created=Tuesday, August 14, 2012 13:45:33}
%TCIDATA{LastRevised=Thursday, February 01, 2024 17:39:50}
%TCIDATA{<META NAME="GraphicsSave" CONTENT="32">}
%TCIDATA{<META NAME="DocumentShell" CONTENT="Exams and Syllabi\SW\Assignment">}

\setlength{\topmargin}{-1.0in}
\setlength{\textheight}{9.25in}
\setlength{\oddsidemargin}{0.0in}
\setlength{\evensidemargin}{0.0in}
\setlength{\textwidth}{6.5in}
\def\labelenumi{\arabic{enumi}.}
\def\theenumi{\arabic{enumi}}
\def\labelenumii{(\alph{enumii})}
\def\theenumii{\alph{enumii}}
\def\p@enumii{\theenumi.}
\def\labelenumiii{\arabic{enumiii}.}
\def\theenumiii{\arabic{enumiii}}
\def\p@enumiii{(\theenumi)(\theenumii)}
\def\labelenumiv{\arabic{enumiv}.}
\def\theenumiv{\arabic{enumiv}}
\def\p@enumiv{\p@enumiii.\theenumiii}
\pagestyle{plain}
\setcounter{secnumdepth}{0}
\parindent=0pt
\input{tcilatex}
\begin{document}


\begin{center}
\begin{equation*}
%TCIMACRO{%
%\FRAME{itbpF}{1.3214in}{0.9496in}{0in}{}{}{Figure}{%
%\special{language "Scientific Word";type "GRAPHIC";maintain-aspect-ratio TRUE;display "USEDEF";valid_file "T";width 1.3214in;height 0.9496in;depth 0in;original-width 1.3552in;original-height 0.9669in;cropleft "0";croptop "1";cropright "1";cropbottom "0";tempfilename 'QYQC7S00.bmp';tempfile-properties "XPR";}}}%
%BeginExpansion
%EndExpansion
\end{equation*}%
\textbf{Electromagnetismo - 230045}

\emph{Gu\'{\i}a N}$%
%TCIMACRO{\U{b0}}%
%BeginExpansion
{{}^\circ}%
%EndExpansion
3$\emph{\ - Distribuciones continuas de carga}

\textit{Profesor: Arturo Fern\'{a}ndez P\'{e}rez}
\end{center}

\section{Distribuciones continuas de carga}

\begin{enumerate}
\item Una barra delgada de longitud $1,69$ [m] tiene una densidad de carga
uniforme $\lambda =4$ $\left[ \frac{mC}{m}\right] $. Determine la carga de,

\begin{enumerate}
\item Un trozo de 70 [cm] de la barra. \textbf{R: }$Q=2,8$ $[mC]$

\item La barra completa. \textbf{R:} $Q=6,76$ $[mC]$
\end{enumerate}

\item Una l\'{\i}nea de transmici\'{o}n tiene una densidad de carga lineal
dada por $\lambda (x)=4\sqrt{x}$ $\left[ \frac{mC}{m}\right] .$ Determine la
carga de un segmento de 100 [m] de longitud (desde $x=0$ [m] a $x=100$ [m]). 
\textbf{R: }$Q=2,666$ $[C]$

\item Un anillo de radio $R=20$ [cm] tiene una densidad lineal de carga
uniforme $\lambda =-5$ $\left[ \frac{nC}{m}\right] .$ Determine la carga de,

\begin{enumerate}
\item Un arco del anillo que subtiende un \'{a}ngulo $\theta =60%
%TCIMACRO{\U{b0}}%
%BeginExpansion
{{}^\circ}%
%EndExpansion
.$ \textbf{R: }$Q=-1,047$ $[nC]$

\item El anillo completo. \textbf{R: }$Q=-6,283$ $[nC]$
\end{enumerate}

\item Considere un cuadrado de lado $a=1,2$ [m] . Determine la carga total
del cuadrado si

\begin{enumerate}
\item \'{E}ste tiene una densidad de carga uniforme $\sigma =-3.5$ $\left[ 
\frac{\mu C}{m^{2}}\right] $ \textbf{R: }$Q=-5,04$ $[\mu C]$
\end{enumerate}

\item Un disco de radio $R=45$ [cm] tiene una distribuci\'{o}n de carga
variable $\sigma =4e^{-r/R}$ $\left[ \frac{\mu C}{m^{2}}\right] $, donde $r$
es la distancia desde el centro. Determine la carga total de disco. \textbf{%
R:} $Q=1,3448$ [$\mu C]$

\item Una esfera macisa de radio $R=50$ [cm] tiene una densidad volum\'{e}%
trica uniforme $\rho =8r^{2}$ $\left[ \frac{\mu C}{m^{3}}\right] $ donde $r$
es la distancia al centro de la esfera. Determine su carga total. \textbf{R:}
$Q=0,628$ [$\mu C]$
\end{enumerate}

\end{document}
